%\ifodd\value{page}\hbox{}\newpage\fi
\chapter*{Śrī Lalitā Sahasranāma — Śloka 1}
\addcontentsline{toc}{chapter}{Śloka 1 - Nomi 1,2,3}

\begin{sloka}
{\sanskritfont
\begin{tabular}{c}
श्रीमाता श्रीमहाराज्ञी श्रीमत्सिंहासनेश्वरी ।\\
चिदग्निकुण्डसम्भूता देवकार्यसमुद्यता ॥ १ ॥
\end{tabular}

\vspace{1.5em}

{\middlesize \iastfont
śrīmātā śrīmahārājñī śrīmatsiṃhāsaneśvarī |\\
cidagnikuṇḍasambhūtā devakāryasamudyatā ||1||\\}

\vspace{1em}

{\middlesize
\anvaya{%
श्रीमाता श्रीमत्सिंहासनेश्वरी श्रीमहाराज्ञी चिदग्निकुण्डसम्भूता देवकार्यसमुद्यता ।}
}

\middlesize \iastfont \itmean{%
«La Madre Suprema, la Grande Sovrana, la Signora del glorioso trono; sorta dal fuoco della coscienza, pronta all’opera per la missione degli dèi.»%
}

}
\end{sloka}

\wordmeaning{%
\wordline{श्री-माता}{śrī-mātā}
  { la Madre Suprema;}%
\wordline{श्री-महा-राज्ञी}{śrī-mahā-rājñī}
  { la Grande Sovrana;}%
\wordline{श्रीमत्-सिंहासन-ईश्वरी}{śrīmat-siṃhāsana-īśvarī}
  { la Signora del glorioso trono;}%
\wordline{चित्-अग्नि-कुण्ड-सम्भूता}{cit-agni-kuṇḍa-sambhūtā}
  { sorta dal fuoco della coscienza;}%
\wordline{देव-कार्य-समुद्यता}{deva-kārya-samudyatā}
  { pronta all’opera per la missione degli dèi;}%
}

Lo \textit{Śloka} 1 introduce \textit{Lalitā} come principio sovrano e originario, stabilendo chi lei è prima di dire che cosa fa. I titoli \textit{Śrīmātā} e \textit{Śrīmahārājñī} affermano la sua natura di Madre assoluta e Regina suprema, mentre \textit{Śrīmatsiṃhāsaneśvarī} la colloca sul trono cosmico, simbolo di autorità ontologica, non politica. La nascita dal \textit{cid-agni-kuṇḍa} indica che la Dea emerge dal fuoco della coscienza, non da una causa materiale: è la coscienza stessa che si rende operativa. \textit{Devakārya-samudyatā} mostra che questa supremazia non è inattiva: \textit{Lalitā} è già in movimento, pronta a compiere l’opera cosmica per l’ordine divino. Questa \textit{śloka} fonda l’intero \textit{Sahasranāma} dichiarando che l’azione cosmica nasce da una sovranità cosciente e materna.

Lo \textit{Śloka} 1 introduce tre \textit{nāma} fondamentali, che stabiliscono l’identità sovrana di \textit{Lalitā} prima di ogni azione. I nomi sono \textbf{\textit{Śrīmātā}}, che la definisce come Madre originaria e principio generativo; \textbf{\textit{Śrīmahārājñī}}, che ne afferma la regalità assoluta e non derivata; e \textbf{\textit{Śrīmatsiṃhāsaneśvarī}}, che la colloca sul trono cosmico come fonte dell’ordine ontologico. Le espressioni \textit{cidagnikuṇḍasambhūtā} e \textit{devakāryasamudyatā} non introducono nuovi nomi, ma qualificano l’origine e la disposizione all’azione di questi tre \textit{nāma}, mostrando che la sovranità nasce dalla coscienza e si traduce immediatamente in operatività cosmica.